% unit: noether.p09.sec.005
\section*{\S{} 5. Les domaines les plus généraux formés d'entiers transcendants, sur la base d'une base algébrique.}

% unit: noether.p09.par.041
Soit $H$ une base algébrique quelconque de tous les nombres, et soit $\mathfrak M(\mathfrak G)$ l'ensemble de tous les domaines $\mathfrak G$ qui contiennent $H$ et donc $[H]$. Lorsque $H$ parcourt toutes les bases possibles, $\mathfrak M(\mathfrak G)$ parcourt l'ensemble de tous les domaines $\mathfrak G$; par conséquent, comme au \S{} 4, nous pouvons nous limiter à considérer les domaines $\mathfrak G$ appartenant à $\mathfrak M(\mathfrak G)$.

% unit: noether.p09.par.042
Les domaines $\mathfrak G$ de $\mathfrak M(\mathfrak G)$ contiennent tous $[H]$ comme diviseur; mais, puisque $[H]$ n'est pas lui-même un domaine $\mathfrak G$, on ne peut pas conclure comme au \S{} 4 que $[H]$ soit le plus grand diviseur commun. Que cela soit néanmoins le cas est montré par une famille dénombrable de domaines $\mathfrak G$ de $\mathfrak M(\mathfrak G)$, obtenue par une légère généralisation du domaine de modules du \S{} 3, et n'ayant en effet aucun élément commun en dehors de $[H]$.

% unit: noether.p09.par.043
\emph{Pour chaque $\nu$ fixé, soient les domaines $\mathfrak G_\nu$ $(\nu=1,2,\ldots\text{ à l'infini})$ composés de toutes les quantités}
% unit: noether.p09.eq.018
\begin{equation}
z=f(\eta)+g_1(\eta)\eta_1^\nu+g_2(\eta)\eta_2^\nu+\cdots+g_t(\eta)\eta_t^\nu,
\tag{1}
\end{equation}
\emph{où $f(\eta)$ parcourt tous les polynômes entiers en les $\eta$; où $\eta_1,\ldots,\eta_t$ parcourent les quantités de base, et où, pour chaque combinaison fixée $\eta_1^\nu,\ldots,\eta_t^\nu$, les fonctions $g_1(\eta),\ldots,g_t(\eta)$ parcourent indépendamment toutes les fonctions algébriques entières des $\eta$.}

% unit: noether.p09.par.044
Comme au \S{} 3, on voit que les domaines $\mathfrak G_\nu$ possèdent les propriétés abstraites I--IV; $\mathfrak G_\nu$ ne contient que des fonctions algébriquement entières et, pour $\nu=1$, il est identique au domaine du \S{} 3. Nous montrons maintenant que les domaines $\mathfrak G_\nu$ n'ont pas d'autre élément commun que les polynômes $f(\eta)$ de $[H]$. Soit en effet $z$ une fonction algébriquement entière quelconque, mais non rationnelle, des $\eta$, satisfaisant à l'équation irréductible relativement à $[H]$
% unit: noether.p09.eq.019
\begin{equation}
z^\chi+a_1(\eta)z^{\chi-1}+a_2(\eta)z^{\chi-2}+\cdots+a_\chi(\eta)=0,
\tag{2}
\end{equation}
où $\chi>1$; et soit $\lambda$ le plus grand des degrés des polynômes $a_1(\eta),\ldots,a_\chi(\eta)$. Si $z$ appartient à $\mathfrak G_\nu$, alors la quantité $\zeta=z-f(\eta)$ satisfait à l'équation, également irréductible relativement à $[H]$,
% unit: noether.p09.eq.020
\begin{equation}
\zeta^\chi+b_1(\eta)\zeta^{\chi-1}+b_2(\eta)\zeta^{\chi-2}+\cdots+b_\chi(\eta)=0,
\tag{3}
\end{equation}
où, d'après (1), $b_1(\eta),b_2(\eta),\ldots,b_\chi(\eta)$, comme fonctions symétriques des $\zeta$, contiennent comme termes de plus basse dimension des termes d'au moins les dimensions $\nu,2\nu,\ldots,\chi\nu$. La comparaison de (2) et (3) donne
% unit: noether.p09.eq.021
\[
 b_1(\eta)=a_1(\eta)+\chi f(\eta).
\]
Si donc la quantité $\xi=z+\frac{a_1(\eta)}{\chi}$ satisfait à
% unit: noether.p09.eq.022
\begin{equation}
\xi^\chi+c_2(\eta)\xi^{\chi-2}+\cdots+c_\chi(\eta)=0,
\tag{4}
\end{equation}
alors
% unit: noether.p09.eq.023
\begin{equation}
 \xi-\zeta=\frac{b_1(\eta)}{\chi};
 \qquad
 c_i(\eta)\equiv b_i(\eta)\pmod{\frac{b_1(\eta)}{\chi}},
\tag{5}
\end{equation}
où $b_1(\eta)$ commence par des termes d'au moins la dimension $\nu$. Or les $c_i(\eta)$, comme le montre la comparaison de (2) avec (4), sont de degré $\chi$ dans les coefficients $a_i(\eta)$ de (2), donc de degré au plus $\chi\cdot\lambda$ dans les $\eta$; d'autre part tous les $b_i(\eta)$ commencent par des termes d'au moins la dimension $\nu$. Ainsi, si l'on choisit $\nu>\chi\cdot\lambda$, (5) donne
% unit: noether.p09.eq.024
\[
 c_i(\eta)=0;\qquad
 \xi^\chi=0\quad\text{ou}\quad
 \left(z+\frac{a_1(\eta)}{\chi}\right)^\chi=0.
\]
Ainsi $z$, contrairement à l'hypothèse, devient rationnelle en les $\eta$. Donc :

% unit: noether.p09.par.045
\emph{Si $z$ est une fonction algébriquement entière, non rationnelle, des $\eta$, alors $z$ ne peut appartenir à aucun domaine $\mathfrak G_\nu$ pour lequel $\nu>\chi\lambda$, ou encore :}

% unit: noether.p09.par.046
\emph{L'intersection de tous les domaines $\mathfrak G$ de $\mathfrak M(\mathfrak G)$ est donnée par le domaine entier $[H]$, qui lui-même n'est pas un domaine $\mathfrak G$; par conséquent $\mathfrak M(\mathfrak G)$ n'est pas fermé pour la formation des intersections.}

% unit: noether.p09.par.047
La construction des domaines $\mathfrak G$ les plus généraux au moyen d'une base algébrique est donc elle aussi moins facile à embrasser que celle des domaines $\mathfrak H$. Le domaine $[H]$ possède les propriétés I, III et IV; si l'on adjoint polynomialement un système quelconque $\mathfrak S$, alors, pour le domaine entier résultant $[H,\mathfrak S]$, les propriétés I et III restent valables, tandis que IV constitue une condition spéciale imposée à $\mathfrak S$. Pour qu'un domaine $\mathfrak G$ naisse, il faut encore imposer II à $\mathfrak S$ comme condition; ou, exprimé autrement : pour tout corps $\mathfrak K$ fini sur $(H)$\footnote{$(H)$ désigne le corps formé de toutes les fonctions rationnelles des $\eta$ à coefficients algébriques.}, il doit exister un corps $\mathfrak L$ au-dessus de $\mathfrak K$, fini sur $\mathfrak K$, tel que $[H,\mathfrak S]$ contienne un élément primitif de $\mathfrak L$.\footnote{Cf. \S{} 7.} Réciproquement, tout domaine $\mathfrak G$ de $\mathfrak M(\mathfrak G)$ peut aussi être représenté sous la forme $[H,\mathfrak S]$; on obtient donc tous les domaines $\mathfrak G$ de $\mathfrak M(\mathfrak G)$ lorsque $\mathfrak S$ parcourt tous les systèmes ainsi restreints. La structure des systèmes $\mathfrak S$ apparaîtra plus simplement dans le paragraphe suivant au moyen de la \og{} base rationnelle \fg{}; en même temps on obtient une analogie complète avec les théorèmes du \S{} 4.

% unit: noether.p09.sec.006
\section*{\S{} 6. Les domaines les plus généraux formés d'entiers transcendants, sur la base d'une base rationnelle.}

% unit: noether.p09.par.048
Par analogie avec la base algébrique, la base rationnelle doit être définie comme suit : \og{} Un système $\Theta$ de nombres $\vartheta$ est appelé base rationnelle de tous les nombres si $\Theta$ permet d'exprimer rationnellement tout nombre --- avec coefficients dans $K$\footnote{Par fonctions rationnelles il faut toujours entendre de telles fonctions à coefficients algébriques; ces coefficients algébriques sont inclus dans la définition à cause de III et IV. Il serait plus adéquat d'exiger des coefficients rationnels aussi bien dans III et IV que dans la définition de la base rationnelle. Les raisonnements des \S{}\S{} 6 et 7 restent littéralement inchangés si l'on remplace seulement le corps $K$ de tous les nombres algébriques par le corps de tous les nombres rationnels.} --- tandis que, sous au moins un bon ordre, aucun nombre de base ne peut être exprimé rationnellement --- avec coefficients dans $K$ --- par un nombre fini de nombres de base précédents. \fg{}\footnote{À la différence des bases linéaire et algébrique, l'ordre des éléments joue ici un rôle. Par exemple, dans l'ordre $\eta,\sqrt[\nu]{\eta}$ les deux éléments doivent être admis, tandis que dans l'ordre inverse seul $\sqrt[\nu]{\eta}$ doit être admis.}

% unit: noether.p09.par.049
Soumettons à un bon ordre $\Omega$ un domaine $\mathfrak G$ formé d'entiers transcendants. Il peut alors arriver qu'une quantité $z$ de $\mathfrak G$ soit exprimable rationnellement par un nombre fini de quantités précédentes de $\mathfrak G$; c'est-à-dire qu'elle satisfasse à une équation $z=\varphi(\xi)$, où $\varphi(\xi)$ est une fonction rationnelle, à coefficients dans $K$, de $\xi_1,\ldots,\xi_t$. En particulier, tout nombre algébrique $\alpha$ du domaine satisfait à l'équation $z=\alpha$. Les quantités restantes $\vartheta$ de $\mathfrak G$ forment un système $\Theta$ dont il faut montrer qu'il est une base rationnelle de toutes les quantités de $\mathfrak G$. En effet, par hypothèse, chaque quantité $\vartheta$ est rationnellement indépendante des quantités précédentes de $\Theta$. L'argument d'induction\footnote{Cf. Zermelo, loc. cit., \S{} 1.} montre en outre que toute relation $z=\varphi(\xi)$ entraîne une relation $z=\psi(\vartheta)$; sinon il devrait y avoir un premier élément $z_0=\varphi_0(\xi_1,\ldots,\xi_t)$ qui ne puisse pas être exprimé rationnellement par les $\vartheta$, tandis que les $\xi_1,\ldots,\xi_t$ précédents seraient des fonctions rationnelles des $\vartheta$; cela donne une contradiction. Par II, $\Theta$ est en même temps une base rationnelle de tous les nombres; par III et I, outre $\Theta$, $\mathfrak G$ contient comme diviseur le domaine entier $[\Theta]$ constitué de tous les polynômes entiers en les $\vartheta$, tandis que, par IV, $[\Theta]$ ne contient aucun nombre algébriquement fractionnaire. Ainsi la base rationnelle $\Theta$ de tous les nombres satisfait à la condition accessoire que le domaine entier $[\Theta]$ qui en provient ne contienne aucun nombre algébriquement fractionnaire;\footnote{Que ce soit réellement une condition est montré par exemple par le fait que les deux quantités $\eta,\frac1{2\sqrt{\eta}}$ ne sont pas admissibles comme éléments de base, tandis que $\eta,\frac1{\sqrt{\eta}}$ le sont.} $\Theta$ est une \emph{base rationnelle avec condition accessoire} pour tous les nombres. Donc, par analogie avec le \S{} 2 : \emph{Tout domaine $\mathfrak G$ peut être construit au moyen d'une base rationnelle avec condition accessoire, de manière que $\mathfrak G$ contienne le domaine entier $[\Theta]$ comme diviseur.}

% unit: noether.p09.par.050
Soit $\Theta$ une base rationnelle quelconque avec condition accessoire pour tous les nombres --- dont l'existence est donc garantie par l'existence des domaines $\mathfrak G$. Soit $\mathfrak N(\mathfrak G)$ l'ensemble de tous les domaines $\mathfrak G$ qui contiennent $\Theta$ et donc $[\Theta]$. Lorsque $\Theta$ parcourt toutes les bases possibles avec condition accessoire, $\mathfrak N(\mathfrak G)$ parcourt, d'après ce qui précède, l'ensemble de tous les domaines $\mathfrak G$; par conséquent nous pouvons à nouveau nous limiter aux domaines $\mathfrak G$ de $\mathfrak N(\mathfrak G)$.

% unit: noether.p09.par.051
Or $[\Theta]$ lui-même est un domaine $\mathfrak G$; car tout nombre peut être exprimé rationnellement par $\Theta$, donc comme quotient de deux polynômes de $[\Theta]$; et, par sa construction, $[\Theta]$ satisfait aussi aux autres conditions I, III, IV. \emph{Ainsi $[\Theta]$ est le plus grand diviseur commun, ou l'intersection, de tous les domaines $\mathfrak G$ de $\mathfrak N(\mathfrak G)$; en outre $\mathfrak N(\mathfrak G)$ est fermé par rapport à l'intersection, c'est-à-dire que l'intersection de tous les domaines $\mathfrak G$ d'un sous-système quelconque de $\mathfrak N(\mathfrak G)$ appartient encore à $\mathfrak N(\mathfrak G)$.} En effet, I vaut pour l'intersection; II et III résultent du fait que $[\Theta]$ est un diviseur; et IV du fait que le domaine d'intersection est contenu comme diviseur dans un domaine $\mathfrak G$.

% unit: noether.p09.par.052
L'adjonction rationnellement entière d'un système quelconque $\mathfrak S$ de fonctions rationnelles des $\vartheta$ à $[\Theta]$ --- toute fonction algébrique des $\vartheta$ pouvant, par la propriété de la base rationnelle, être exprimée rationnellement au moyen de certains autres $\vartheta$ --- donne un domaine entier $[\Theta,\mathfrak S]$ possédant les propriétés I, II, III, et devient donc un domaine $\mathfrak G$ dès que $\mathfrak S$ est un \og{} système admissible \fg{}, c'est-à-dire dès qu'il satisfait à la condition qu'aucun nombre algébriquement fractionnaire n'entre dans le domaine par l'adjonction rationnellement entière. Réciproquement, tout domaine $\mathfrak G$ de $\mathfrak N(\mathfrak G)$ permet la représentation $[\Theta,\mathfrak S]$, pourvu seulement que $\mathfrak S$ désigne le système résiduel $\mathfrak G-[\Theta]$. Ainsi, pour les domaines $\mathfrak G$ les plus généraux, en complète analogie avec le \S{} 4, on a :
\begin{enumerate}
\item[a)] \emph{L'intersection de tous les domaines $\mathfrak G$ de $\mathfrak N(\mathfrak G)$ est donnée par le domaine entier $[\Theta]$, qui est lui-même un domaine $\mathfrak G$; $\mathfrak N(\mathfrak G)$ est fermé par rapport à la formation des intersections.}
\item[b)] \emph{Tout domaine $\mathfrak G$ de $\mathfrak N(\mathfrak G)$ naît par adjonction rationnellement entière d'un système admissible $\mathfrak S$ à $[\Theta]$. Lorsque $\Theta$ parcourt toutes les bases rationnelles possibles avec condition accessoire, $\mathfrak N(\mathfrak G)$ parcourt l'ensemble de tous les domaines $\mathfrak G$.}\footnote{Pour les systèmes admissibles $\mathfrak S$ les plus généraux, ce qui a été dit au \S{} 4 vaut en remplaçant \og{} algébrique \fg{} par \og{} rationnel \fg{}.}
\end{enumerate}

% unit: noether.p09.par.053
\emph{Remarque.} Si l'on extrait de $\Theta$ une base algébrique $H$ de toutes les quantités de $\Theta$, alors $H$ est en même temps une base algébrique de tous les nombres; réciproquement toute base algébrique peut être prolongée en une base rationnelle en plaçant $H$ d'abord dans le bon ordre. En conséquence, la construction du \S{} 5 des domaines $\mathfrak G$ au moyen d'une base algébrique peut s'interpréter comme suit. Décomposons le système $\mathfrak S$ à adjoindre à $[H]$ en deux sous-systèmes $\mathfrak S_1$ et $\mathfrak S_2$; l'adjonction de $\mathfrak S_1$ revient à prolonger $H$ en une base rationnelle avec condition accessoire, à laquelle on adjoint ensuite encore polynomialement un système admissible $\mathfrak S_2$.

% unit: noether.p09.sec.007
\section*{\S{} 7. Construction de la base rationnelle la plus générale avec condition accessoire.}

% unit: noether.p09.par.054
Les résultats du paragraphe précédent ont encore besoin d'un complément, car on y a prouvé l'\emph{existence} d'une base rationnelle avec condition accessoire pour tous les nombres, mais non un procédé pour la construire à partir d'un bon ordre prescrit. À cause de la condition accessoire, cette construction prend une forme plus compliquée lorsque, au lieu d'un domaine $\mathfrak G$, on a devant soi le corps de tous les nombres complexes.\footnote{Au \S{} 6 b), il suffirait de ne faire parcourir à $\Theta$ que les bases rationnelles spéciales qui se composent d'une base algébrique et de fonctions algébriquement entières des quantités de base; la condition accessoire est alors automatiquement satisfaite. Pour le voir, il suffit, d'après \S{} 6, de prolonger en une base rationnelle une base algébrique extraite de $\mathfrak G$, puis, dans la base ainsi obtenue, de multiplier chaque fonction algébriquement fractionnaire des $\eta$ par un polynôme convenablement choisi en les $\eta$, ce qui conserve la propriété de base. Cette construction se transpose sans changement au corps de tous les nombres complexes. La question, issue du \S{} 6, de la base rationnelle la plus générale avec condition accessoire possède toutefois aussi un intérêt propre.}

% unit: noether.p09.par.055
Soumettons le corps de tous les nombres complexes à un bon ordre $\Omega$. Alors des quantités de trois espèces peuvent se présenter :
\begin{enumerate}
\item[1.] Des quantités $y$ dont l'adjonction rationnellement entière aux quantités de base précédentes --- définies récursivement en 3.\footnote{Comme le montre l'argument d'induction, une telle définition récursive est possible.} --- force l'entrée d'un nombre algébriquement fractionnaire dans le domaine entier ainsi obtenu; si aucune quantité de base ne les précède, ce domaine consiste en tous les polynômes entiers en $y$. En particulier, tous les nombres algébriquement fractionnaires sont des quantités $y$, puisque le domaine entier obtenu contient les quantités $y=\beta$.
\item[2.] Des quantités $z$ qui ne possèdent pas la propriété 1, mais qui peuvent être exprimées rationnellement par un nombre fini de nombres précédents différents des $y$; elles satisfont donc à une relation $z=\varphi(\xi_1,\ldots,\xi_t)$, où $\varphi$ est une fonction rationnelle à coefficients dans $K$ et où aucun des $\xi$ n'est un $y$. En particulier, tous les nombres algébriques entiers $\alpha$ appartiennent ici, puisque la propriété 1 ne vaut pas pour eux et qu'ils satisfont à la relation $z=\alpha$.
\item[3.] Des quantités $\vartheta$ pour lesquelles ni 1 ni 2 ne vaut et qui doivent être appelées quantités de base. En particulier, le premier nombre transcendant $\vartheta_0$ dans le bon ordre est une telle quantité de base; car les polynômes entiers en $\vartheta_0$ ne peuvent représenter aucun nombre algébriquement fractionnaire, et $\vartheta_0$ ne peut pas non plus dépendre rationnellement des nombres algébriques qui le précèdent.
\end{enumerate}

% unit: noether.p09.par.056
\emph{Il faut maintenant montrer que le système $\Theta$ de toutes les quantités de base $\vartheta$ est une base rationnelle avec condition accessoire de tous les nombres.}

% unit: noether.p09.par.057
Comme aucune quantité $y$ ne se trouve parmi les $\vartheta$, le domaine entier $[\Theta]$ ne peut contenir aucun nombre algébriquement fractionnaire; la condition accessoire est donc satisfaite. Comme, en outre, aucune quantité $z$ ne se trouve parmi les $\vartheta$, chaque $\vartheta$ est rationnellement indépendante des quantités de base précédentes. L'argument d'induction montre, exactement comme au \S{} 6, que toute relation $z=\varphi(\xi)$ --- où aucun $y$ ne se trouve parmi les $\xi$ --- entraîne une relation $z=\psi(\vartheta)$, de sorte que toutes les quantités $z$ sont exprimables rationnellement par les $\vartheta$.

% unit: noether.p09.par.058
Il reste seulement à prouver l'assertion plus difficile que les $y$ sont eux aussi exprimables rationnellement par les $\vartheta$. Nous montrons d'abord que les $y$ dépendent algébriquement des $\vartheta$. Par hypothèse, il existe pour chaque $y$ au moins un nombre algébriquement fractionnaire $\beta$ admettant la représentation
% unit: noether.p09.eq.025
\begin{equation}
\beta=f_0(\vartheta)+f_1(\vartheta)y+\cdots+f_\chi(\vartheta)y^\chi,
\tag{1}
\end{equation}
où $f_0(\vartheta),\ldots,f_\chi(\vartheta)$ sont des polynômes de $[\Theta]$ et ne peuvent donc représenter aucun nombre algébriquement fractionnaire. Si $y$ était algébriquement indépendante des $\vartheta$, (1) vaudrait identiquement en $y$ et l'on aurait $\beta=f_0(\vartheta)$, en contradiction avec les propriétés de $[\Theta]$.

% unit: noether.p09.par.059
Pour conclure de la dépendance algébrique à la dépendance rationnelle, nous extrayons de $\Theta$ une base algébrique $H$ de toutes les quantités de $\Theta$. Alors chaque $y$, comme fonction algébrique des $\vartheta$, dépend aussi algébriquement des $\eta$; par $y=y(\eta)$ est donc déterminé un corps algébrique $(H,y(\eta))$ au-dessus de $(H)$. Comme tous les $\eta$ sont en même temps des quantités $\vartheta$, la preuve de la dépendance rationnelle revient à montrer que pour chacun de ces corps il existe des éléments primitifs qui sont des fonctions rationnelles des $\vartheta$; plus précisément, on montrera que $y(\eta)$ perd la propriété 1 après multiplication par un polynôme en les $\eta$, et passe par conséquent en un tel élément.

% unit: noether.p09.par.060
On sait qu'entre $(H)$ et $(H,y)$ il n'existe qu'un nombre fini de corps intermédiaires. Soient $\Omega_0=(H),\Omega_1,\ldots,\Omega_\sigma$ ceux d'entre eux qui possèdent un élément primitif exprimable rationnellement par les $\vartheta$, de sorte que tout élément du corps devient une fonction rationnelle des $\vartheta$; cette condition est au moins toujours satisfaite pour $\Omega_0=(H)$. Soit donnée la fonction primitive appartenant à $\Omega_i$ sous la forme $\frac{g_i(\vartheta)}{h_i(\vartheta)}$, où $g_i,h_i$ sont des polynômes de $[\Theta]$. Alors --- $\alpha_i$ désignant un entier convenablement choisi --- le polynôme de $[\Theta]$
% unit: noether.p09.eq.026
\[
 \xi_i=g_i(\vartheta)+\alpha_i h_i(\vartheta)
\]
est une fonction primitive de $\Omega_i$, ou d'un corps algébrique au-dessus de $\Omega_i$,\footnote{Cf. la fin du \S{} 5.} et toute fonction $\psi(\eta)$ de $\Omega_i$ admet la représentation
% unit: noether.p09.eq.027
\begin{equation}
\psi(\eta)=
\frac{a_1(\eta)\xi_i^{\kappa-1}+a_2(\eta)\xi_i^{\kappa-2}+\cdots+a_\kappa(\eta)}
     {a(\eta)}
=\frac{A(\eta,\xi_i)}{a(\eta)},
\tag{2}
\end{equation}
où $a(\eta),a_1(\eta),\ldots,a_\kappa(\eta)$ sont des polynômes entiers en les $\eta$, et donc $A(\eta,\xi_i)$ et $a(\eta)$ sont des quantités de $[\Theta]$.

% unit: noether.p09.par.061
L'équation irréductible satisfaite par $y$ relativement à $\Omega_i$ est, d'après (2), de la forme
% unit: noether.p09.eq.028
\[
 f^{(i)}(\eta)y^{\lambda_i}+f_1^{(i)}(\eta,\xi_i)y^{\lambda_i-1}
+\cdots+f_{\lambda_i}^{(i)}(\eta,\xi_i)=0,
\]
et tous les coefficients de cette équation appartiennent à $[\Theta]$. Alors la fonction $y_i=f^{(i)}(\eta)\cdot y$ satisfait à l'équation également irréductible, relativement à $\Omega_i$,
% unit: noether.p09.eq.029
\[
 y_i^{\lambda_i}+f_1^{(i)}(\eta,\xi_i)y_i^{\lambda_i-1}
+\cdots+\bigl(f^{(i)}(\eta)\bigr)^{\lambda_i-1}
 f_{\lambda_i}^{(i)}(\eta,\xi_i)=0,
\]
en vertu de laquelle elle est entière sur $[H,\xi_i]$; il en est de même du produit de $y_i$ par n'importe quel polynôme en les $\eta$. En particulier, la fonction
% unit: noether.p09.eq.030
\begin{equation}
Y=f^{(0)}(\eta)f^{(1)}(\eta)\cdots f^{(\sigma)}(\eta)\cdot y
\tag{3}
\end{equation}
est entière sur chacun des domaines $[H,\xi_i]$ en vertu de l'équation, irréductible dans chaque cas relativement à $\Omega_i$,
% unit: noether.p09.eq.031
\begin{equation}
Y^{\lambda_i}+F_1^{(i)}(\eta,\xi_i)Y^{\lambda_i-1}
+\cdots+F_{\lambda_i}^{(i)}(\eta,\xi_i)=0
\qquad (i=0,1,\ldots,\sigma).
\tag{4}
\end{equation}
Nous montrons maintenant que $Y$ est l'élément primitif recherché, exprimable rationnellement par les $\vartheta$. Supposons par exemple que le nombre algébrique $\gamma$ soit représenté par $Y$ et $[\Theta]$ :
% unit: noether.p09.eq.032
\begin{equation}
\gamma=k_0(\vartheta)+k_1(\vartheta)Y+\cdots+k_\nu(\vartheta)Y^\nu,
\tag{5}
\end{equation}
où $k_0(\vartheta),\ldots,k_\nu(\vartheta)$ sont des polynômes de $[\Theta]$. Déterminons maintenant l'équation irréductible satisfaite par $Y$ relativement au corps $\mathfrak K=(H;k_0(\vartheta),\ldots,k_\nu(\vartheta))$ provenant du domaine de coefficients de (5). Elle est, comme on sait, donnée par l'équation irréductible de $Y$ relativement au corps d'intersection de $\mathfrak K$ et $(H,Y)$, situé entre $(H)$ et $(H,Y)$. Or $(H,Y)$ est identique à $(H,y)$ d'après (3); et comme tous les éléments de $\mathfrak K$, donc aussi tous les éléments du corps d'intersection, peuvent être exprimés rationnellement par un nombre fini de $\vartheta$, ce dernier corps est nécessairement l'un des corps $\Omega_0,\Omega_1,\ldots,\Omega_\sigma$, disons $\Omega_\tau$. L'équation irréductible cherchée est donc, d'après (4), de degré $\lambda_\tau$ et s'écrit
% unit: noether.p09.eq.033
\begin{equation}
Y^{\lambda_\tau}+F_1^{(\tau)}(\eta,\xi_\tau)Y^{\lambda_\tau-1}
+\cdots+F_{\lambda_\tau}^{(\tau)}(\eta,\xi_\tau)=0
\tag{6}
\end{equation}
avec des polynômes de $[\Theta]$ comme coefficients. Au moyen de (6), (5) peut être réécrite sous la forme
% unit: noether.p09.eq.034
\begin{equation}
\gamma=K_0(\vartheta)+K_1(\vartheta)Y+\cdots+
K_{\lambda_\tau-1}(\vartheta)Y^{\lambda_\tau-1},
\tag{7}
\end{equation}
où $K_0(\vartheta),K_1(\vartheta),\ldots$ appartiennent à $\mathfrak K$ et, d'autre part, sont des polynômes de $[\Theta]$. Mais puisque $\gamma$ est un nombre algébrique, la relation (7) n'atteint pas le degré $\lambda_\tau$ en $Y$ et est donc satisfaite \emph{identiquement} en $Y$. Ainsi
% unit: noether.p09.eq.035
\[
 \gamma=K_0(\vartheta),
\]
c'est-à-dire que $\gamma$ appartient à $[\Theta]$ et est par conséquent un entier algébrique.

% unit: noether.p09.par.062
Lorsque $\gamma$ parcourt tous les nombres algébriques représentables par $Y$ et $[\Theta]$, le corps d'intersection de $(H,Y)$ et du corps correspondant $\mathfrak K$ ne parcourt toujours que les corps $\Omega_0,\Omega_1,\ldots,\Omega_\sigma$, et tous les arguments précédents restent valables. Autrement dit, \emph{par adjonction rationnellement entière de $Y$ à $[\Theta]$, aucun nombre algébriquement fractionnaire ne peut être représenté}. Ainsi $Y$ ne possède plus la propriété 1; selon 2) ou 3) c'est une quantité $z$ ou $\vartheta$, et donc elle est rationnellement représentable par un nombre fini de $\vartheta$; d'après (3), il en est de même de $y$ : \emph{$\Theta$ est bien une base rationnelle de tous les nombres.} Réciproquement, puisque toute base rationnelle avec condition accessoire peut être construite suivant 1), 2), 3) en plaçant $\Theta$ d'abord dans le bon ordre, nous avons prouvé : \emph{La construction donnée par 1), 2), 3) conduit à la base rationnelle la plus générale avec condition accessoire de tous les nombres.}

% unit: noether.p09.sec.008
\section*{\S{} 8. Répartition des domaines $\mathfrak G$ et $\mathfrak H$ en classes.\srcfn{*)}{Voir E. Steinitz, \emph{Algebraische Theorie der Körper}, Journal de Crelle 137 (1910), en particulier les \S\S{} 23 et 24. Le \S{} 23 établit déjà l'existence de la base algébrique.}}

% unit: noether.p09.par.063
À côté de la question des domaines $\mathfrak G$ et $\mathfrak H$ les plus généraux se pose la question supplémentaire des domaines \emph{essentiellement différents}, qui --- dans un sens à préciser ci-dessous --- ne peuvent pas être engendrés par le même principe de construction. Cela conduit à une répartition de ces domaines en classes.

% unit: noether.p09.par.064
\emph{Deux domaines seront dits appartenir à la même classe, ou être isomorphes, si leurs éléments peuvent être mis en correspondance biunivoque de telle sorte que la somme et le produit de deux éléments quelconques de l'un soient toujours associés à la somme et au produit des éléments correspondants de l'autre.} Il résulte de cette définition que, dans toute relation isomorphe, zéro et l'unité correspondent à eux-mêmes, et que toutes les relations algébriques entre un nombre fini d'éléments sont conservées pour les éléments correspondants, et réciproquement. En particulier, outre l'unité, tous les nombres rationnels correspondent à eux-mêmes; tandis que la totalité des nombres algébriques --- et de même la totalité des entiers algébriques --- se transforme en elle-même, bien qu'un nombre algébrique particulier puisse très bien correspondre à un conjugué.

% unit: noether.p09.par.065
Remarquons d'abord que tout domaine de l'ensemble $\mathfrak M(\mathfrak G)$ de tous les domaines $\mathfrak G$ contenant une base algébrique fixe $H$ est isomorphe à un domaine de l'ensemble $\mathfrak M^*(\mathfrak G)$ appartenant à une autre base algébrique $H^*$. En effet, le corps de tous les nombres complexes se déduit de chaque base algébrique par extension algébrique et a donc la même cardinalité que la base;\footnote{Steinitz, \S{}\S{} 23 et 24.} ainsi $H$ et $H^*$ ont la même cardinalité, et l'application isomorphe cherchée s'obtient en faisant correspondre les quantités de base $\eta_i$ aux $\eta_i^*$.\footnote{Cette conclusion ne vaut naturellement pas pour les ensembles $\mathfrak N(\mathfrak G)$ et $\mathfrak N^*(\mathfrak G)$ appartenant à deux bases rationnelles $\Theta$ et $\Theta^*$; l'isomorphisme de $\mathfrak N(\mathfrak G)$ et de $\mathfrak N^*(\mathfrak G)$ ne peut pas non plus être déduit de l'exprimabilité rationnelle mutuelle de $\Theta$ et $\Theta^*$. Ce que l'on obtient en revanche, c'est une relation isomorphe entre les corps $(\Theta)$ et $(\Theta^*)$.} D'autre part, $\mathfrak M(\mathfrak G)$ contient des domaines $\mathfrak G$ non isomorphes, ou essentiellement différents, comme il résulte des exemples des \S{}\S{} 2, 3 et 5. Car, puisque toutes les relations algébriques sont conservées dans l'application isomorphe, une base algébrique doit correspondre à nouveau à une base algébrique; or les \S{}\S{} 2 et 3 ont montré que les domaines de fonctionnelles et de modules ne possèdent aucune base algébrique pour laquelle toutes les propriétés de base du domaine de Zermelo seraient satisfaites. Comme le domaine de fonctionnelles du \S{} 2 est intégralement clos, le sous-ensemble $\mathfrak M(\mathfrak H)$ contient lui aussi des domaines $\mathfrak H$ essentiellement différents.

% unit: noether.p09.par.066
Nous montrons maintenant, de façon analogue à la construction de bases, comment on peut extraire de $\mathfrak M(\mathfrak G)$ un sous-ensemble $\mathfrak L(\mathfrak G)$ tel que tous les domaines $\mathfrak G$ de $\mathfrak L(\mathfrak G)$ soient essentiellement différents, tandis que tout domaine de $\mathfrak M(\mathfrak G)$ --- et donc tout domaine $\mathfrak G$ en général --- soit isomorphe à un domaine de $\mathfrak L(\mathfrak G)$. Soumettons $\mathfrak M(\mathfrak G)$ à un bon ordre $\Omega$; alors un domaine $\mathfrak G$ de $\mathfrak M(\mathfrak G)$ peut être isomorphe ou non à un domaine précédent. Les domaines $\mathfrak G$ isomorphes à aucun précédent forment un ensemble $\mathfrak L(\mathfrak G)$, qui, par construction, ne contient que des domaines $\mathfrak G$ essentiellement différents. L'argument d'induction montre aussi que tout domaine de $\mathfrak M(\mathfrak G)$ est isomorphe à un domaine de $\mathfrak L(\mathfrak G)$. Car si $\mathfrak G_1$ était le premier domaine pour lequel cela échouerait, alors $\mathfrak G_1$ appartiendrait lui-même à $\mathfrak L(\mathfrak G)$, ou bien serait isomorphe à un domaine précédent $\mathfrak G_0$, lequel, par hypothèse, est isomorphe à un domaine de $\mathfrak L(\mathfrak G)$; il en irait donc de même de $\mathfrak G_1$. Puisque tout domaine $\mathfrak G$ appartient à un certain ensemble $\mathfrak M^*(\mathfrak G)$ et est donc isomorphe à un domaine de $\mathfrak M(\mathfrak G)$, \emph{la totalité des classes de domaines $\mathfrak G$ est représentée par $\mathfrak L(\mathfrak G)$.}

% unit: noether.p09.par.067
Si, dans un domaine $\mathfrak G$ de $\mathfrak L(\mathfrak G)$, on remplace les quantités de base $\eta$ par celles $\eta^*$ d'une autre base, alors, comme on l'a montré plus haut, il en résulte un domaine isomorphe à $\mathfrak G$. Réciproquement, soit $\mathfrak G^*$ un domaine quelconque et supposons-le isomorphe à un domaine $\mathfrak G$ de $\mathfrak L(\mathfrak G)$. Si les quantités de base $\eta$ correspondent aux quantités $\eta^*$ de $\mathfrak G^*$, celles-ci forment à nouveau une base algébrique, et $\mathfrak G^*$ contient les mêmes fonctions algébriques de $\eta^*$ que $\mathfrak G$ contient de $\eta$ --- ou leurs conjuguées. \emph{Ainsi tout domaine isomorphe à $\mathfrak G$ s'obtient en faisant parcourir à la base algébrique $H$ toutes les bases algébriques possibles dans $\mathfrak G$ et dans ses conjugués relativement à $(H)$}\footnote{Les éléments de ces domaines conjugués à $\mathfrak G$ peuvent certes être exprimés rationnellement par les éléments de $\mathfrak G$ d'après II, mais pas nécessairement intégralement et rationnellement; les conjugués peuvent donc différer de $\mathfrak G$.} \emph{--- si ceux-ci diffèrent de $\mathfrak G$.} À partir de chaque domaine fixé $\mathfrak G_i$ dans $\mathfrak L(\mathfrak G)$, on obtient ainsi la classe $\mathfrak R_i$, qui contient tous et seulement les domaines $\mathfrak G$ isomorphes à $\mathfrak G_i$, bien que chaque domaine puisse y apparaître plusieurs fois, voire une infinité de fois. De $\mathfrak R_i$ on peut de nouveau extraire, par le procédé précédent, un sous-ensemble $\mathfrak R_i^*$ qui contient chaque domaine isomorphe à $\mathfrak G_i$ une et une seule fois. Lorsque $\mathfrak R_i^*$ parcourt toutes les classes, on obtient la totalité de tous les domaines $\mathfrak G$, chacun une et une seule fois. Pour caractériser la classe $\mathfrak R_i$, on peut naturellement employer, au lieu de $\mathfrak G_i$, tout autre domaine de la classe dont tous les domaines de $\mathfrak R_i$ peuvent être déduits comme ci-dessus. En résumé :

% unit: noether.p09.par.068
\emph{De $\mathfrak M(\mathfrak G)$ on peut extraire un sous-ensemble $\mathfrak L(\mathfrak G)$ tel que les domaines $\mathfrak G$ de $\mathfrak L(\mathfrak G)$ soient en correspondance biunivoque avec la totalité des classes (sans répétition). À partir de $\mathfrak G_i$, la classe correspondante $\mathfrak R_i$ naît en faisant parcourir à $H$ toutes les bases algébriques possibles dans $\mathfrak G_i$ et ses conjugués relativement à $(H)$. Si $\mathfrak R_i^*$ désigne tous les domaines mutuellement distincts de $\mathfrak R_i$, et si $\mathfrak R_i^*$ parcourt toutes les classes, alors la totalité de tous les domaines $\mathfrak G$ apparaît, chaque domaine une et une seule fois.}\footnote{Ici $\mathfrak M(\mathfrak G)$ doit être construit conformément aux \S{}\S{} 6 et 7 : d'après \S{} 7, on prolonge une base algébrique $H$ en chaque base rationnelle $\Theta$ avec condition accessoire qui en provient, et, pour chaque $\Theta$, on forme d'après \S{} 6 b) la totalité $\mathfrak N(\mathfrak G)$ de tous les domaines $\mathfrak G$ contenant $\Theta$.}

% unit: noether.p09.par.069
L'énoncé analogue vaut pour les domaines $\mathfrak H$ formés d'entiers transcendants algébriquement entiers.

% unit: noether.p09.sec.009
\section*{\S{} 9. Un exemple d'une infinité dénombrable de classes de domaines $\mathfrak G$.}

% unit: noether.p09.par.070
La cardinalité des constructions données au \S{} 4 b) et au \S{} 6 b) se lit, d'après les remarques qui les accompagnent, comme égale à la cardinalité de tous les bons ordres, donc comme $2^c$, si $c$ désigne la cardinalité du continu. Mais on ne peut rien en conclure sur la cardinalité de la totalité de tous les domaines $\mathfrak G$, chacun compté une seule fois, et encore moins sur la cardinalité de toutes les classes. Que le nombre de classes ne soit certainement pas fini sera maintenant montré par un exemple d'une infinité dénombrable de domaines $\mathfrak G$ --- analogues à ceux du \S{} 5 --- qui se révéleront tous essentiellement différents et fourniront ainsi une infinité dénombrable de classes.\footnote{Les domaines $\mathfrak G$ du \S{} 6 donnent eux aussi une infinité dénombrable de classes, mais la preuve est un peu plus compliquée.}

% unit: noether.p09.par.071
Par analogie avec \S{} 5, (1), soit le domaine $\mathfrak G_\sigma$ formé de la totalité des quantités
% unit: noether.p09.eq.036
\begin{equation}
 z\equiv a_0+a_1(\eta)+\cdots+a_{\sigma-1}(\eta)
 \pmod{\mathfrak M_\sigma(\eta)},
\tag{1}
\end{equation}
\emph{où $a_i(\eta)$ désigne des formes homogènes de $i$-ième dimension en les $\eta$, et où le module $\mathfrak M_\sigma$ contient pour base tous les produits de puissances de dimension $\sigma$ en les $\eta$}\footnote{Naturellement, dans chaque $z$ individuel, seuls un nombre fini de produits de puissances des $\eta$ apparaissent dans le module.} \emph{et pour multiplicateurs toutes les fonctions algébriquement entières des $\eta$.} Nous montrons que, sous une relation isomorphe entre deux domaines $\mathfrak G_\sigma$ et $\mathfrak G_\tau$ $(\sigma\ne\tau)$, les modules $\mathfrak M_\sigma$ et $\mathfrak M_\tau$ doivent nécessairement se correspondre isomorphiquement; l'impossibilité s'ensuit alors facilement.

% unit: noether.p09.par.072
Désignons par $\xi$ les indéterminées de $\mathfrak G_\tau$; sous la relation isomorphe, supposons que les quantités $f(\eta)$ de $\mathfrak G_\sigma$ correspondent à ces $\xi$. Comme la relation isomorphe se conserve lors de la formation des quotients, le domaine $\mathfrak G_\xi$ de toutes les fonctions algébriquement entières des $\xi$ correspond à un domaine intégralement clos $\mathfrak T$ formé de toutes les fonctions entières sur les $f(\eta)$, et donc aussi entières sur le domaine engendré par les $\eta$.

% unit: noether.p09.par.073
Supposons que (1) corresponde à une quantité de $\mathfrak M_\tau$. Alors, par la propriété de module de $\mathfrak M_\tau$, le produit de (1) par une quantité arbitraire $\psi(\eta)$ de $\mathfrak T$ doit lui aussi appartenir à $\mathfrak G_\sigma$; en formules :
% unit: noether.p09.eq.037
\begin{equation}
(a_0+a_1(\eta)+\cdots+a_{\sigma-1}(\eta))\psi(\eta)
\equiv b_0+b_1(\eta)+\cdots+b_{\sigma-1}(\eta)
\pmod{\mathfrak M_\sigma}.
\tag{2}
\end{equation}
En annulant tous les $\eta$, on obtient $a_0\psi(0)=b_0$, donc (2) devient
% unit: noether.p09.eq.038
\begin{equation}
(a_0+a_1(\eta)+\cdots+a_{\sigma-1}(\eta))(\psi(\eta)-\psi(0))
\equiv c_1(\eta)+\cdots+c_{\sigma-1}(\eta)
\pmod{\mathfrak M_\sigma}.
\tag{3}
\end{equation}
Or, outre $\psi(\eta)$, les fonctions
% unit: noether.p09.eq.039
\[
 \chi_\nu(\eta)=\sqrt[\nu]{\psi(\eta)-\psi(0)}
\]
appartiennent aussi à $\mathfrak T$ pour tout $\nu$; et, puisque $\chi_\nu(0)=0$, les formules analogues à (3) sont :
% unit: noether.p09.eq.040
\begin{equation}
(a_0+a_1(\eta)+\cdots+a_{\sigma-1}(\eta))\chi_\nu(\eta)
\equiv c_1^{(\nu)}(\eta)+\cdots+c_{\sigma-1}^{(\nu)}(\eta)
\pmod{\mathfrak M_\sigma}
\quad(\nu=1,2,\ldots).
\tag{4}
\end{equation}
Soit $A(\eta)$, de degré $\lambda$, le produit de $\chi_1=\psi(\eta)-\psi(0)$ par ses $\kappa-1$ conjugués. Comme $\chi_1(0)=0$, $A(\eta)$ doit commencer par des termes d'au moins la première dimension. De (4) on tire
% unit: noether.p09.eq.041
\begin{equation}
a_0\cdot\chi_\nu(\eta)\equiv0\pmod{\mathfrak M_1};\qquad
a_0^\nu\cdot\chi_1(\eta)\equiv0\pmod{\mathfrak M_\nu};\qquad
a_0^{\nu\kappa}\cdot A(\eta)\equiv0\pmod{\mathfrak M_{\nu\kappa}},
\tag{5}
\end{equation}
et donc, pour $a_0\ne0$, comme au \S{} 3, $\lambda>\nu\kappa$. Mais comme $\nu<\frac{\lambda}{\kappa}$ contredit la clôture algébrique-entière de $\mathfrak T$, il faut nécessairement $a_0=0$. De (4) on obtient alors
% unit: noether.p09.eq.042
\begin{equation}
a_1(\eta)\cdot\chi_\nu(\eta)\equiv c_1^{(\nu)}(\eta)\pmod{\mathfrak M_2};\qquad
[a_1(\eta)]^{\nu\kappa}\cdot A(\eta)
\equiv[c_1^{(\nu)}(\eta)]^{\nu\kappa}\pmod{\mathfrak M_{\nu\kappa+1}},
\tag{6}
\end{equation}
et de là, puisque $A(\eta)$ commence par des termes d'au moins la première dimension, $c_1^{(\nu)}(\eta)=0$. Ainsi, en complète analogie avec (5), pour tout $\nu$ :
% unit: noether.p09.eq.043
\[
 a_1(\eta)\cdot\chi_\nu(\eta)\equiv0\pmod{\mathfrak M_2};\qquad
 [a_1(\eta)]^{\nu\kappa}\cdot A(\eta)\equiv0\pmod{\mathfrak M_{2\nu\kappa}},
\]
et donc $a_1(\eta)=0$. En continuant de cette manière, l'application alternée de (5) et (6) donne
% unit: noether.p09.eq.044
\[
 a_0=0,\qquad a_1(\eta)=0,\quad \ldots,\quad a_{\sigma-1}(\eta)=0.
\]
Comme la considération analogue vaut aussi relativement à $\mathfrak G_\tau$ : \emph{Sous toute relation isomorphe entre $\mathfrak G_\sigma$ et $\mathfrak G_\tau$ $(\sigma\ne\tau)$, les modules $\mathfrak M_\sigma$ et $\mathfrak M_\tau$ se correspondent isomorphiquement.}

% unit: noether.p09.par.074
Supposons par exemple que $\sigma>\tau$. Aux indéterminées $\xi$ correspondaient les quantités $f(\eta)$ de $\mathfrak G_\sigma$; et, puisque toutes les quantités de $\mathfrak G_\tau$ peuvent être représentées par des polynômes en les $\xi$ mod. $\mathfrak M_\tau$, la correspondance de $\mathfrak M_\tau$ et de $\mathfrak M_\sigma$ implique que toutes les quantités de $\mathfrak G_\sigma$ s'obtiennent comme polynômes en les $f(\eta)$ mod. $\mathfrak M_\sigma$. Comme $\sigma>\tau\ge1$, les indéterminées $\eta$ n'appartiennent certainement pas à $\mathfrak M_\sigma$, et donc, puisqu'elles sont exprimables intégralement et rationnellement par les $f(\eta)$ mod. $\mathfrak M_\sigma$, il doit exister un $f(\eta)$ ayant des termes linéaires non nuls. Soit donc
% unit: noether.p09.eq.045
\[
 \xi:f(\eta)\equiv a_0+a_1(\eta)\pmod{\mathfrak M_2}\quad[a_1(\eta)\ne0],
\]
% unit: noether.p09.eq.046
\[
 \xi^\tau:[f(\eta)]^\tau\equiv a_0^\tau\pmod{\mathfrak M_1}\quad\text{pour }a_0\ne0,
\]
% unit: noether.p09.eq.047
\[
 \hphantom{\xi^\tau:[f(\eta)]^\tau}\equiv [a_1(\eta)]^\tau
 \pmod{\mathfrak M_{\tau+1}}\quad\text{pour }a_0=0.
\]
Or $\xi^\tau$ appartient à $\mathfrak M_\tau$, tandis que $[f(\eta)]^\tau$ commence par des termes non nuls de dimension zéro ou de dimension $\tau$ et, puisque $\tau<\sigma$, ne peut appartenir à $\mathfrak M_\sigma$, contrairement à la correspondance de $\mathfrak M_\sigma$ et $\mathfrak M_\tau$ obtenue ci-dessus. Le cas $\tau>\sigma$ est traité en même temps en échangeant $\eta$ et $\xi$. Il est donc montré que les domaines $\mathfrak G_\sigma$ et $\mathfrak G_\tau$ sont \emph{non isomorphes, ou essentiellement différents}. Ainsi \emph{les domaines $\mathfrak G_i$ $(i=1,2,\ldots\text{ à l'infini})$ donnent un exemple d'une infinité dénombrable de classes de domaines $\mathfrak G$.}

% unit: noether.p09.par.075
\emph{Remarque.} Bien que deux domaines quelconques $\mathfrak G_i$ soient non isomorphes, il peut très bien arriver, pour deux domaines, que chacun soit isomorphe à une partie de l'autre.\footnote{Ce comportement peut aussi se produire pour les corps; cf. Steinitz, loc. cit., conclusion.} Soient par exemple $\tau=1$ et $\sigma$ arbitraire :
% unit: noether.p09.eq.048
\[
\mathfrak G_1:a_0+\mathfrak M_1(\xi);\qquad
\mathfrak G_\sigma:a_0+a_1(\eta)+\cdots+a_{\sigma-1}(\eta)+\mathfrak M_\sigma(\eta).
\]
L'application $\xi=\eta$ fait de $\mathfrak G_\sigma$ un diviseur propre de $\mathfrak G_1$; inversement, la correspondance biunivoque dans $\mathfrak G_1$ et $\mathfrak G_\sigma$ donnée par $\xi=\eta^\sigma$, $\eta=\sqrt[\sigma]{\xi}$ fait de $\mathfrak G_1$ un diviseur propre de $\mathfrak G_\sigma$. Ainsi chaque $\mathfrak G_i$ est aussi isomorphe à un diviseur propre de lui-même. À la différence de la notion d'intersection, la notion de classe d'intersection --- c'est-à-dire d'une classe telle que chacun de ses domaines serait isomorphe à un diviseur propre de chaque domaine de chaque autre classe --- n'existe pas, du moins pas comme notion déterminée de manière unique.

% unit: noether.p09.sec.010
\section*{\S{} 10. Les quantités entières d'un corps arbitraire.}

% unit: noether.p09.par.076
Les développements des paragraphes précédents ont reposé, pour les domaines $\mathfrak G$ --- pourvu que les nombres algébriques ne soient pas inclus dans les propriétés de définition III et IV des domaines $\mathfrak G$ ---, seulement sur le fait que la totalité de tous les nombres complexes forme un corps;\footnote{Ce n'est qu'au \S{} 7 qu'a été utilisée l'hypothèse supplémentaire que le \og{} corps premier \fg{} des nombres rationnels contient une infinité d'éléments; dans le cas d'un corps premier ne contenant qu'un nombre fini d'éléments, ce paragraphe disparaît simplement, comme nous allons le voir.} pour les domaines $\mathfrak H$, ils ont reposé en outre sur le fait que ce corps est algébriquement clos. Les résultats obtenus admettent donc directement la \emph{généralisation au corps le plus général}, qui, d'après Weber et E. Steinitz,\footnote{loc. cit., introduction.} est défini abstraitement comme \og{} un système d'éléments muni de deux opérations, addition et multiplication, soumises aux lois associative et commutative, liées par la loi distributive, et admettant des opérations inverses sans restriction et uniques \fg{}.\footnote{Seule la division par zéro est à exclure.}

% unit: noether.p09.par.077
Tout corps de ce genre contient un élément unité, et par conséquent le \og{} corps premier \fg{} qui en dérive par les opérations d'addition et de multiplication et leurs inverses. Ce corps premier est soit du type des nombres rationnels (caractéristique $0$), soit du type du système des classes modulo un nombre premier $p$ (caractéristique $p$).\footnote{Steinitz, \S{} 4.} Les \emph{quantités entières du corps premier} sont alors bien définies comme les quantités issues de l'élément unité par addition et soustraction; pour les corps premiers de caractéristique $p$, les quantités entières épuisent déjà le corps premier et il n'existe pas de quantités fractionnaires du corps premier. Tout corps algébriquement clos contient un \og{} corps absolument algébrique \fg{}, obtenu à partir du corps premier par adjonction de tous les éléments algébriques sur le corps premier; pour les corps de caractéristique $0$, il est donc du type du corps de tous les nombres algébriques. Les \emph{quantités algébriquement entières du corps absolument algébrique} sont alors bien définies comme toutes les quantités algébriquement entières sur les quantités entières du corps premier (ou, ce qui revient au même, sur l'unité); pour les corps de caractéristique $p$, il n'existe donc pas non plus de quantités algébriquement fractionnaires du corps absolument algébrique. Après ces remarques préliminaires, la définition des quantités \og{} entières \fg{} et \og{} algébriquement entières \fg{} se transpose aux corps arbitraires, respectivement aux corps algébriquement clos : \emph{Les quantités entières d'un corps $\Omega$ sont définies comme un domaine entier $\mathfrak G$ de quantités de $\Omega$ dont le corps des fractions\footnote{C'est-à-dire le corps formé de tous les quotients de deux quantités de $\mathfrak G$.} épuise $\Omega$, et qui contient l'élément unité mais aucune quantité fractionnaire du corps premier.}

% unit: noether.p09.par.078
\emph{Les quantités entières au sens algébrique d'un corps algébriquement clos $\Omega$ sont définies comme un domaine entier $\mathfrak H$ intégralement clos de quantités de $\Omega$ dont le corps des fractions épuise $\Omega$, et qui contient l'élément unité mais aucune quantité algébriquement fractionnaire du corps absolument algébrique.}

% unit: noether.p09.par.079
Pour les corps de caractéristique zéro, les résultats obtenus dans les paragraphes précédents s'appliquent littéralement dès que l'on remplace seulement les \og{} nombres algébriques \fg{} par les quantités du corps premier, respectivement du corps absolument algébrique. Pour les corps de caractéristique $p$, les résultats se simplifient encore, parce que le corps premier ne contient aucune quantité fractionnaire, de sorte que la condition accessoire pour la base rationnelle et pour les systèmes à adjoindre disparaît simplement.\footnote{Cf. la première note de ce paragraphe.} Ainsi le résultat est le suivant : \emph{La définition abstraite permet comme quantités entières d'un corps de caractéristique $p$ : tout domaine entier dérivé d'une base rationnelle quelconque, le corps lui-même, et tout domaine entier situé entre ces domaines et le corps. De façon correspondante, comme quantités entières au sens algébrique d'un corps algébriquement clos de caractéristique première, on a : tout domaine entier intégralement clos situé entre le domaine dérivé d'une base algébrique et le corps lui-même --- y compris les deux domaines limites.} Pour les corps les plus généraux de caractéristique première, comme pour les corps premiers correspondants, la distinction entre entier et fractionnaire s'efface donc.

% unit: noether.p09.close.001
\medskip
\noindent Erlangen, 30 mars 1915.
